% Options for packages loaded elsewhere
\PassOptionsToPackage{unicode}{hyperref}
\PassOptionsToPackage{hyphens}{url}
\documentclass[
]{article}
\usepackage{xcolor}
\usepackage[top=2cm, bottom=2cm, left=2cm, right=2cm]{geometry}
\usepackage{amsmath,amssymb}
\setcounter{secnumdepth}{5}
\usepackage{iftex}
\ifPDFTeX
  \usepackage[T1]{fontenc}
  \usepackage[utf8]{inputenc}
  \usepackage{textcomp} % provide euro and other symbols
\else % if luatex or xetex
  \usepackage{unicode-math} % this also loads fontspec
  \defaultfontfeatures{Scale=MatchLowercase}
  \defaultfontfeatures[\rmfamily]{Ligatures=TeX,Scale=1}
\fi
\usepackage{lmodern}
\ifPDFTeX\else
  % xetex/luatex font selection
\fi
% Use upquote if available, for straight quotes in verbatim environments
\IfFileExists{upquote.sty}{\usepackage{upquote}}{}
\IfFileExists{microtype.sty}{% use microtype if available
  \usepackage[]{microtype}
  \UseMicrotypeSet[protrusion]{basicmath} % disable protrusion for tt fonts
}{}
\makeatletter
\@ifundefined{KOMAClassName}{% if non-KOMA class
  \IfFileExists{parskip.sty}{%
    \usepackage{parskip}
  }{% else
    \setlength{\parindent}{0pt}
    \setlength{\parskip}{6pt plus 2pt minus 1pt}}
}{% if KOMA class
  \KOMAoptions{parskip=half}}
\makeatother
\usepackage{color}
\usepackage{fancyvrb}
\newcommand{\VerbBar}{|}
\newcommand{\VERB}{\Verb[commandchars=\\\{\}]}
\DefineVerbatimEnvironment{Highlighting}{Verbatim}{commandchars=\\\{\}}
% Add ',fontsize=\small' for more characters per line
\usepackage{framed}
\definecolor{shadecolor}{RGB}{248,248,248}
\newenvironment{Shaded}{\begin{snugshade}}{\end{snugshade}}
\newcommand{\AlertTok}[1]{\textcolor[rgb]{0.94,0.16,0.16}{#1}}
\newcommand{\AnnotationTok}[1]{\textcolor[rgb]{0.56,0.35,0.01}{\textbf{\textit{#1}}}}
\newcommand{\AttributeTok}[1]{\textcolor[rgb]{0.13,0.29,0.53}{#1}}
\newcommand{\BaseNTok}[1]{\textcolor[rgb]{0.00,0.00,0.81}{#1}}
\newcommand{\BuiltInTok}[1]{#1}
\newcommand{\CharTok}[1]{\textcolor[rgb]{0.31,0.60,0.02}{#1}}
\newcommand{\CommentTok}[1]{\textcolor[rgb]{0.56,0.35,0.01}{\textit{#1}}}
\newcommand{\CommentVarTok}[1]{\textcolor[rgb]{0.56,0.35,0.01}{\textbf{\textit{#1}}}}
\newcommand{\ConstantTok}[1]{\textcolor[rgb]{0.56,0.35,0.01}{#1}}
\newcommand{\ControlFlowTok}[1]{\textcolor[rgb]{0.13,0.29,0.53}{\textbf{#1}}}
\newcommand{\DataTypeTok}[1]{\textcolor[rgb]{0.13,0.29,0.53}{#1}}
\newcommand{\DecValTok}[1]{\textcolor[rgb]{0.00,0.00,0.81}{#1}}
\newcommand{\DocumentationTok}[1]{\textcolor[rgb]{0.56,0.35,0.01}{\textbf{\textit{#1}}}}
\newcommand{\ErrorTok}[1]{\textcolor[rgb]{0.64,0.00,0.00}{\textbf{#1}}}
\newcommand{\ExtensionTok}[1]{#1}
\newcommand{\FloatTok}[1]{\textcolor[rgb]{0.00,0.00,0.81}{#1}}
\newcommand{\FunctionTok}[1]{\textcolor[rgb]{0.13,0.29,0.53}{\textbf{#1}}}
\newcommand{\ImportTok}[1]{#1}
\newcommand{\InformationTok}[1]{\textcolor[rgb]{0.56,0.35,0.01}{\textbf{\textit{#1}}}}
\newcommand{\KeywordTok}[1]{\textcolor[rgb]{0.13,0.29,0.53}{\textbf{#1}}}
\newcommand{\NormalTok}[1]{#1}
\newcommand{\OperatorTok}[1]{\textcolor[rgb]{0.81,0.36,0.00}{\textbf{#1}}}
\newcommand{\OtherTok}[1]{\textcolor[rgb]{0.56,0.35,0.01}{#1}}
\newcommand{\PreprocessorTok}[1]{\textcolor[rgb]{0.56,0.35,0.01}{\textit{#1}}}
\newcommand{\RegionMarkerTok}[1]{#1}
\newcommand{\SpecialCharTok}[1]{\textcolor[rgb]{0.81,0.36,0.00}{\textbf{#1}}}
\newcommand{\SpecialStringTok}[1]{\textcolor[rgb]{0.31,0.60,0.02}{#1}}
\newcommand{\StringTok}[1]{\textcolor[rgb]{0.31,0.60,0.02}{#1}}
\newcommand{\VariableTok}[1]{\textcolor[rgb]{0.00,0.00,0.00}{#1}}
\newcommand{\VerbatimStringTok}[1]{\textcolor[rgb]{0.31,0.60,0.02}{#1}}
\newcommand{\WarningTok}[1]{\textcolor[rgb]{0.56,0.35,0.01}{\textbf{\textit{#1}}}}
\usepackage{graphicx}
\makeatletter
\newsavebox\pandoc@box
\newcommand*\pandocbounded[1]{% scales image to fit in text height/width
  \sbox\pandoc@box{#1}%
  \Gscale@div\@tempa{\textheight}{\dimexpr\ht\pandoc@box+\dp\pandoc@box\relax}%
  \Gscale@div\@tempb{\linewidth}{\wd\pandoc@box}%
  \ifdim\@tempb\p@<\@tempa\p@\let\@tempa\@tempb\fi% select the smaller of both
  \ifdim\@tempa\p@<\p@\scalebox{\@tempa}{\usebox\pandoc@box}%
  \else\usebox{\pandoc@box}%
  \fi%
}
% Set default figure placement to htbp
\def\fps@figure{htbp}
\makeatother
\setlength{\emergencystretch}{3em} % prevent overfull lines
\providecommand{\tightlist}{%
  \setlength{\itemsep}{0pt}\setlength{\parskip}{0pt}}
\usepackage{titling}
\setlength{\droptitle}{-1.5cm}
\pretitle{\begin{center}\Large\bfseries}
\posttitle{\end{center}}
\usepackage{bookmark}
\IfFileExists{xurl.sty}{\usepackage{xurl}}{} % add URL line breaks if available
\urlstyle{same}
\hypersetup{
  pdftitle={PHÂN TÍCH THỐNG KÊ NHIỀU CHIỀU (MAT3452)},
  pdfauthor={Bài tập tuần 6: Phương pháp làm giảm số chiều dữ liệu},
  hidelinks,
  pdfcreator={LaTeX via pandoc}}

\title{\textbf{PHÂN TÍCH THỐNG KÊ NHIỀU CHIỀU (MAT3452)}}
\author{\textbf{Bài tập tuần 6: Phương pháp làm giảm số chiều dữ liệu}}
\date{Trần Đăng Quang Minh - MSV: 23000141}

\begin{document}
\maketitle

\begin{center}\rule{0.5\linewidth}{0.5pt}\end{center}

\textbf{Câu 1. } Tập dữ liệu \textbf{USJudgeRatings {datasets}} trong R
gồm 12 chỉ số đánh giá cho 43 thẩm phán. Sử dụng dữ liệu trên, hãy cho
biết:

\hfill\break
(1) Thực hiện phân tích thành phần chính trên tập dữ liệu, hãy cho biết
nếu dùng 5 thành phần chính đầu thì tỉ lệ dữ liệu là bao nhiêu?

\begin{Shaded}
\begin{Highlighting}[]
\FunctionTok{library}\NormalTok{(datasets)}

\NormalTok{df }\OtherTok{\textless{}{-}}\NormalTok{ USJudgeRatings}
\FunctionTok{head}\NormalTok{(df)}
\end{Highlighting}
\end{Shaded}

\begin{verbatim}
##                CONT INTG DMNR DILG CFMG DECI PREP FAMI ORAL WRIT PHYS RTEN
## AARONSON,L.H.   5.7  7.9  7.7  7.3  7.1  7.4  7.1  7.1  7.1  7.0  8.3  7.8
## ALEXANDER,J.M.  6.8  8.9  8.8  8.5  7.8  8.1  8.0  8.0  7.8  7.9  8.5  8.7
## ARMENTANO,A.J.  7.2  8.1  7.8  7.8  7.5  7.6  7.5  7.5  7.3  7.4  7.9  7.8
## BERDON,R.I.     6.8  8.8  8.5  8.8  8.3  8.5  8.7  8.7  8.4  8.5  8.8  8.7
## BRACKEN,J.J.    7.3  6.4  4.3  6.5  6.0  6.2  5.7  5.7  5.1  5.3  5.5  4.8
## BURNS,E.B.      6.2  8.8  8.7  8.5  7.9  8.0  8.1  8.0  8.0  8.0  8.6  8.6
\end{verbatim}

\begin{Shaded}
\begin{Highlighting}[]
\CommentTok{\# Phân tích thành phần chính}
\NormalTok{pca }\OtherTok{\textless{}{-}} \FunctionTok{princomp}\NormalTok{(df)}
\FunctionTok{summary}\NormalTok{(pca)}
\end{Highlighting}
\end{Shaded}

\begin{verbatim}
## Importance of components:
##                           Comp.1     Comp.2     Comp.3     Comp.4      Comp.5
## Standard deviation     2.9944812 0.97631616 0.54821237 0.46436892 0.260897305
## Proportion of Variance 0.8476102 0.09010189 0.02840865 0.02038352 0.006434164
## Cumulative Proportion  0.8476102 0.93771205 0.96612070 0.98650423 0.992938391
##                             Comp.6      Comp.7      Comp.8       Comp.9
## Standard deviation     0.172098999 0.127444176 0.107411902 0.0862216855
## Proportion of Variance 0.002799688 0.001535299 0.001090581 0.0007027259
## Cumulative Proportion  0.995738079 0.997273378 0.998363959 0.9990666848
##                             Comp.10      Comp.11      Comp.12
## Standard deviation     0.0703843063 0.0550575991 0.0434546479
## Proportion of Variance 0.0004682789 0.0002865415 0.0001784947
## Cumulative Proportion  0.9995349638 0.9998215053 1.0000000000
\end{verbatim}

\begin{Shaded}
\begin{Highlighting}[]
\CommentTok{\# Nếu dùng 5 thành phần chính}
\FunctionTok{sum}\NormalTok{(pca}\SpecialCharTok{$}\NormalTok{sdev[}\DecValTok{1}\SpecialCharTok{:}\DecValTok{5}\NormalTok{]}\SpecialCharTok{\^{}}\DecValTok{2}\NormalTok{) }\SpecialCharTok{/} \FunctionTok{sum}\NormalTok{(pca}\SpecialCharTok{$}\NormalTok{sdev}\SpecialCharTok{\^{}}\DecValTok{2}\NormalTok{)}
\end{Highlighting}
\end{Shaded}

\begin{verbatim}
## [1] 0.9929384
\end{verbatim}

\hfill\break
(2) Gọi 5 thành phần chính là \(Y_1\) đến \(Y_5\). Hãy cho biết biến
``CONT'' và ``CFMG'' đóng góp bao nhiêu trong mỗi thành phần trên.

\begin{Shaded}
\begin{Highlighting}[]
\NormalTok{contribution }\OtherTok{\textless{}{-}} \FunctionTok{loadings}\NormalTok{(pca)[}\FunctionTok{c}\NormalTok{(}\StringTok{"CONT"}\NormalTok{, }\StringTok{"CFMG"}\NormalTok{), }\DecValTok{1}\SpecialCharTok{:}\DecValTok{5}\NormalTok{]}
\FunctionTok{colnames}\NormalTok{(contribution) }\OtherTok{\textless{}{-}} \FunctionTok{c}\NormalTok{(}\StringTok{"Y1"}\NormalTok{, }\StringTok{"Y2"}\NormalTok{, }\StringTok{"Y3"}\NormalTok{, }\StringTok{"Y4"}\NormalTok{, }\StringTok{"Y5"}\NormalTok{)}
\NormalTok{contribution}
\end{Highlighting}
\end{Shaded}

\begin{verbatim}
##               Y1        Y2         Y3          Y4          Y5
## CONT  0.00599117 0.9332488  0.3198540  0.11293271  0.09462326
## CFMG -0.27201855 0.1631993 -0.1893545 -0.02481803 -0.47964024
\end{verbatim}

\hfill\break
(3) Vẽ biểu đồ tỉ lệ các thành phần trong trường hợp phân tích trên tập
dữ liệu và sử dụng ma trận tương quan.

\begin{Shaded}
\begin{Highlighting}[]
\FunctionTok{library}\NormalTok{(showtext)}
\end{Highlighting}
\end{Shaded}

\begin{verbatim}
## Loading required package: sysfonts
\end{verbatim}

\begin{verbatim}
## Loading required package: showtextdb
\end{verbatim}

\begin{Shaded}
\begin{Highlighting}[]
\FunctionTok{font\_add}\NormalTok{(}\StringTok{"LM"}\NormalTok{, }\StringTok{"C:/Windows/Fonts/cambria.ttc"}\NormalTok{)}
\FunctionTok{showtext\_auto}\NormalTok{()}

\NormalTok{pcacov }\OtherTok{\textless{}{-}} \FunctionTok{princomp}\NormalTok{(df, }\AttributeTok{cor =} \ConstantTok{TRUE}\NormalTok{)}
\FunctionTok{summary}\NormalTok{(pcacov, }\AttributeTok{cor =} \ConstantTok{TRUE}\NormalTok{)}
\end{Highlighting}
\end{Shaded}

\begin{verbatim}
## Importance of components:
##                           Comp.1     Comp.2    Comp.3     Comp.4      Comp.5
## Standard deviation     3.1833165 1.05078398 0.5769763 0.50383231 0.290607615
## Proportion of Variance 0.8444586 0.09201225 0.0277418 0.02115392 0.007037732
## Cumulative Proportion  0.8444586 0.93647089 0.9642127 0.98536661 0.992404341
##                             Comp.6      Comp.7      Comp.8       Comp.9
## Standard deviation     0.193095982 0.140295449 0.124158319 0.0885069038
## Proportion of Variance 0.003107172 0.001640234 0.001284607 0.0006527893
## Cumulative Proportion  0.995511513 0.997151747 0.998436354 0.9990891437
##                             Comp.10      Comp.11      Comp.12
## Standard deviation     0.0749114592 0.0570804224 0.0453913429
## Proportion of Variance 0.0004676439 0.0002715146 0.0001716978
## Cumulative Proportion  0.9995567876 0.9998283022 1.0000000000
\end{verbatim}

\begin{Shaded}
\begin{Highlighting}[]
\NormalTok{prop\_var1 }\OtherTok{\textless{}{-}}\NormalTok{ pca}\SpecialCharTok{$}\NormalTok{sdev}\SpecialCharTok{\^{}}\DecValTok{2} \SpecialCharTok{/} \FunctionTok{sum}\NormalTok{(pca}\SpecialCharTok{$}\NormalTok{sdev}\SpecialCharTok{\^{}}\DecValTok{2}\NormalTok{)}
\FunctionTok{plot}\NormalTok{(prop\_var1,}
  \AttributeTok{col =} \StringTok{"red"}\NormalTok{,}
  \AttributeTok{cex =} \DecValTok{2}\NormalTok{,}
  \AttributeTok{lwd =} \DecValTok{2}\NormalTok{,}
  \AttributeTok{type =} \StringTok{"b"}\NormalTok{,}
  \AttributeTok{xlab =} \StringTok{"Thành phần chính"}\NormalTok{,}
  \AttributeTok{ylab =} \StringTok{"Tỉ lệ"}\NormalTok{,}
  \AttributeTok{main =} \StringTok{"Tỉ lệ các thành phần (phân tích trên bộ dữ liệu)"}
\NormalTok{)}
\end{Highlighting}
\end{Shaded}

\pandocbounded{\includegraphics[keepaspectratio]{W6_files/figure-latex/p1.3-1.pdf}}

\begin{Shaded}
\begin{Highlighting}[]
\NormalTok{prop\_var2 }\OtherTok{\textless{}{-}}\NormalTok{ pcacov}\SpecialCharTok{$}\NormalTok{sdev}\SpecialCharTok{\^{}}\DecValTok{2} \SpecialCharTok{/} \FunctionTok{sum}\NormalTok{(pcacov}\SpecialCharTok{$}\NormalTok{sdev}\SpecialCharTok{\^{}}\DecValTok{2}\NormalTok{)}
\FunctionTok{plot}\NormalTok{(prop\_var2,}
  \AttributeTok{col =} \StringTok{"red"}\NormalTok{,}
  \AttributeTok{cex =} \DecValTok{2}\NormalTok{,}
  \AttributeTok{lwd =} \DecValTok{2}\NormalTok{,}
  \AttributeTok{type =} \StringTok{"b"}\NormalTok{,}
  \AttributeTok{xlab =} \StringTok{"Thành phần chính"}\NormalTok{,}
  \AttributeTok{ylab =} \StringTok{"Tỉ lệ"}\NormalTok{,}
  \AttributeTok{main =} \StringTok{"Tỉ lệ các thành phần (dùng ma trận tương quan)"}
\NormalTok{)}
\end{Highlighting}
\end{Shaded}

\pandocbounded{\includegraphics[keepaspectratio]{W6_files/figure-latex/p1.3-2.pdf}}

\hfill\break
(4) Vẽ biểu đồ hai thành phần trong trường hợp sử dụng ma trận tương
quan.

\begin{Shaded}
\begin{Highlighting}[]
\FunctionTok{library}\NormalTok{(showtext)}
\FunctionTok{font\_add}\NormalTok{(}\StringTok{"LM"}\NormalTok{, }\StringTok{"C:/Windows/Fonts/cambria.ttc"}\NormalTok{)}
\FunctionTok{showtext\_auto}\NormalTok{()}

\NormalTok{pcacov }\OtherTok{\textless{}{-}} \FunctionTok{princomp}\NormalTok{(df, }\AttributeTok{cor =} \ConstantTok{TRUE}\NormalTok{)}

\FunctionTok{biplot}\NormalTok{(pcacov,}
  \AttributeTok{col =} \FunctionTok{c}\NormalTok{(}\DecValTok{2}\NormalTok{, }\DecValTok{4}\NormalTok{),}
  \AttributeTok{cex =} \FunctionTok{c}\NormalTok{(.}\DecValTok{5}\NormalTok{, }\FloatTok{1.25}\NormalTok{),}
  \AttributeTok{xlab =} \StringTok{"Thành phần chính thứ nhất"}\NormalTok{,}
  \AttributeTok{ylab =} \StringTok{"Thành phần chính thứ hai"}\NormalTok{,}
  \AttributeTok{main =} \StringTok{"Biểu đồ hai thành phần (sử dụng ma trận tương quan)"}
\NormalTok{)}
\end{Highlighting}
\end{Shaded}

\pandocbounded{\includegraphics[keepaspectratio]{W6_files/figure-latex/p1.4-1.pdf}}

\hfill\break
(5) Tiến hành phân tích nhân tố cho bộ dữ liệu trên với 4 nhân tố. Hãy
cho biết tỉ lệ tổng bình phương sai của 4 nhân tố.

\begin{Shaded}
\begin{Highlighting}[]
\FunctionTok{factanal}\NormalTok{(df, }\AttributeTok{factors =} \DecValTok{4}\NormalTok{)}
\end{Highlighting}
\end{Shaded}

\begin{verbatim}
## 
## Call:
## factanal(x = df, factors = 4)
## 
## Uniquenesses:
##  CONT  INTG  DMNR  DILG  CFMG  DECI  PREP  FAMI  ORAL  WRIT  PHYS  RTEN 
## 0.749 0.013 0.032 0.023 0.014 0.025 0.008 0.005 0.005 0.005 0.047 0.006 
## 
## Loadings:
##      Factor1 Factor2 Factor3 Factor4
## CONT                          0.499 
## INTG  0.726   0.596   0.175  -0.273 
## DMNR  0.666   0.616   0.305  -0.228 
## DILG  0.916   0.345   0.121         
## CFMG  0.874   0.326   0.259   0.220 
## DECI  0.886   0.285   0.267   0.195 
## PREP  0.935   0.270   0.212         
## FAMI  0.944   0.202   0.241         
## ORAL  0.881   0.335   0.323         
## WRIT  0.911   0.298   0.262         
## PHYS  0.682   0.290   0.622   0.127 
## RTEN  0.797   0.458   0.380         
## 
##                Factor1 Factor2 Factor3 Factor4
## SS loadings      7.829   1.650   1.084   0.507
## Proportion Var   0.652   0.138   0.090   0.042
## Cumulative Var   0.652   0.790   0.880   0.922
## 
## Test of the hypothesis that 4 factors are sufficient.
## The chi square statistic is 49.78 on 24 degrees of freedom.
## The p-value is 0.00151
\end{verbatim}

\hfill\break
(6) Với nhiều nhân tố hơn, hãy cho biết cần bao nhiêu nhân tố để tỉ lệ
tổng phương sai lớn hơn 99\%. Khi đó hãy cho biết biến ``CONT'' và
``CFMG'' được viết thông qua các nhân tố \("F_i"\) như thế nào?

\begin{Shaded}
\begin{Highlighting}[]
\NormalTok{fa7 }\OtherTok{\textless{}{-}} \FunctionTok{factanal}\NormalTok{(df, }\AttributeTok{factors =} \DecValTok{7}\NormalTok{)}
\NormalTok{fa7}\SpecialCharTok{$}\NormalTok{loadings}
\end{Highlighting}
\end{Shaded}

\begin{verbatim}
## 
## Loadings:
##      Factor1 Factor2 Factor3 Factor4 Factor5 Factor6 Factor7
## CONT                          0.994                         
## INTG  0.578   0.783   0.143                                 
## DMNR  0.520   0.795   0.262  -0.119                         
## DILG  0.872   0.435   0.159                  -0.142         
## CFMG  0.845   0.381   0.303   0.130                         
## DECI  0.870   0.346   0.293           0.164                 
## PREP  0.863   0.445   0.214                                 
## FAMI  0.862   0.425   0.220          -0.102   0.104         
## ORAL  0.787   0.527   0.304                                 
## WRIT  0.819   0.507   0.236                                 
## PHYS  0.626   0.368   0.674                                 
## RTEN  0.696   0.611   0.367                                 
## 
##                Factor1 Factor2 Factor3 Factor4 Factor5 Factor6 Factor7
## SS loadings      6.493   3.127   1.123   1.040   0.059   0.048   0.016
## Proportion Var   0.541   0.261   0.094   0.087   0.005   0.004   0.001
## Cumulative Var   0.541   0.802   0.895   0.982   0.987   0.991   0.992
\end{verbatim}

\hfill\break
(7) Trong phân tích nhân tố nói trên, hãy cho biết 4 nhân tố đã đủ để
giải thích bộ dữ liệu hay chưa, nếu chưa thì hãy cho biết cần bao nhiêu
nhân tô? Xét ở mức ý nghĩa 5\% cho kiểm định đó.

\begin{Shaded}
\begin{Highlighting}[]
\FunctionTok{factanal}\NormalTok{(df, }\AttributeTok{factors =} \DecValTok{5}\NormalTok{)}
\end{Highlighting}
\end{Shaded}

\begin{verbatim}
## 
## Call:
## factanal(x = df, factors = 5)
## 
## Uniquenesses:
##  CONT  INTG  DMNR  DILG  CFMG  DECI  PREP  FAMI  ORAL  WRIT  PHYS  RTEN 
## 0.693 0.021 0.025 0.005 0.005 0.028 0.007 0.005 0.005 0.005 0.016 0.007 
## 
## Loadings:
##      Factor1 Factor2 Factor3 Factor4 Factor5
## CONT                          0.551         
## INTG  0.687   0.644   0.182  -0.242         
## DMNR  0.615   0.694   0.286  -0.183         
## DILG  0.908   0.336   0.180           0.156 
## CFMG  0.849   0.363   0.291   0.237         
## DECI  0.864   0.320   0.296   0.186         
## PREP  0.912   0.319   0.243                 
## FAMI  0.917   0.279   0.257                 
## ORAL  0.845   0.411   0.330                 
## WRIT  0.878   0.382   0.266                 
## PHYS  0.640   0.313   0.684                 
## RTEN  0.758   0.510   0.392                 
## 
##                Factor1 Factor2 Factor3 Factor4 Factor5
## SS loadings      7.283   2.097   1.246   0.510   0.047
## Proportion Var   0.607   0.175   0.104   0.042   0.004
## Cumulative Var   0.607   0.782   0.885   0.928   0.932
## 
## Test of the hypothesis that 5 factors are sufficient.
## The chi square statistic is 31.85 on 16 degrees of freedom.
## The p-value is 0.0105
\end{verbatim}

\hfill\break
\textbf{Câu 2. } Bộ dữ liệu \textbf{wdbc {kdevine}} trong R được thu
thập từ Bệnh viện Đại học Wisconsin, Madison, và được sử dụng để hỗ trợ
trong việc chẩn đoán ung thư vú thông qua các đặc điểm của tế bào. Bộ dữ
liệu gồm 569 quan sát và 31 biến, cụ thể các đặc trưng: Radius (Độ dài
trung bình từ tâm đến các điểm trên chu vi), Texture (Độ lệch chuẩn của
các giá trị xám), Perimeter (Chu vi của tế bào), Area (Diện tích của tế
bào), Smoothness (Biến thiên cục bộ trong độ dài bán kính), Compactness
(Tính toán từ chu vi và diện tích), Concavity (Mức độ nghiêm trọng của
các phần lõm trên đường viền), Concave Points (Số lượng phần lõm trên
đường viền), Symmetry (Độ đối xứng của tế bào), Fractal Dimension (Đánh
giá độ phức tạp của hình dạng tế bào).

\hfill\break
(1) Tiến hành tiền xử lý dữ liệu, bao gồm việc loại bỏ các cột không cần
thiết và chuyển đổi các biến phân loại thành biến số.

\begin{Shaded}
\begin{Highlighting}[]
\FunctionTok{library}\NormalTok{(kdevine)}
\end{Highlighting}
\end{Shaded}

\begin{verbatim}
## The kdevine package is no longer actively developed. Consider using 
##    - the 'kde1d' package for marginal estimation, 
##    - the functions vine() and vinecop() from the 'rvinecopulib' 
##      package as replacements for kdevine() and kdevinecop().
\end{verbatim}

\begin{Shaded}
\begin{Highlighting}[]
\NormalTok{df }\OtherTok{\textless{}{-}}\NormalTok{ wdbc}

\NormalTok{df}\SpecialCharTok{$}\NormalTok{ID }\OtherTok{\textless{}{-}} \ConstantTok{NULL}

\NormalTok{df[, }\DecValTok{1}\NormalTok{] }\OtherTok{\textless{}{-}} \FunctionTok{as.numeric}\NormalTok{(df[, }\DecValTok{1}\NormalTok{]) }\SpecialCharTok{{-}} \DecValTok{1}

\FunctionTok{head}\NormalTok{(df)}
\end{Highlighting}
\end{Shaded}

\begin{verbatim}
##   diagnosis mean radius mean texture mean perimeter mean area mean smoothness
## 1         1       17.99        10.38         122.80    1001.0         0.11840
## 2         1       20.57        17.77         132.90    1326.0         0.08474
## 3         1       19.69        21.25         130.00    1203.0         0.10960
## 4         1       11.42        20.38          77.58     386.1         0.14250
## 5         1       20.29        14.34         135.10    1297.0         0.10030
## 6         1       12.45        15.70          82.57     477.1         0.12780
##   mean compactness mean concavity mean concave points mean symmetry
## 1          0.27760         0.3001             0.14710        0.2419
## 2          0.07864         0.0869             0.07017        0.1812
## 3          0.15990         0.1974             0.12790        0.2069
## 4          0.28390         0.2414             0.10520        0.2597
## 5          0.13280         0.1980             0.10430        0.1809
## 6          0.17000         0.1578             0.08089        0.2087
##   mean fractal dimension worst radius worst texture worst perimeter worst area
## 1                0.07871       1.0950        0.9053           8.589     153.40
## 2                0.05667       0.5435        0.7339           3.398      74.08
## 3                0.05999       0.7456        0.7869           4.585      94.03
## 4                0.09744       0.4956        1.1560           3.445      27.23
## 5                0.05883       0.7572        0.7813           5.438      94.44
## 6                0.07613       0.3345        0.8902           2.217      27.19
##   worst smoothness worst compactness worst concavity worst concave points
## 1         0.006399           0.04904         0.05373              0.01587
## 2         0.005225           0.01308         0.01860              0.01340
## 3         0.006150           0.04006         0.03832              0.02058
## 4         0.009110           0.07458         0.05661              0.01867
## 5         0.011490           0.02461         0.05688              0.01885
## 6         0.007510           0.03345         0.03672              0.01137
##   worst symmetry worst fractal dimension sd radius sd texture sd perimeter
## 1        0.03003                0.006193     25.38      17.33       184.60
## 2        0.01389                0.003532     24.99      23.41       158.80
## 3        0.02250                0.004571     23.57      25.53       152.50
## 4        0.05963                0.009208     14.91      26.50        98.87
## 5        0.01756                0.005115     22.54      16.67       152.20
## 6        0.02165                0.005082     15.47      23.75       103.40
##   sd area sd smoothness sd compactness sd concavity sd concave points
## 1  2019.0        0.1622         0.6656       0.7119            0.2654
## 2  1956.0        0.1238         0.1866       0.2416            0.1860
## 3  1709.0        0.1444         0.4245       0.4504            0.2430
## 4   567.7        0.2098         0.8663       0.6869            0.2575
## 5  1575.0        0.1374         0.2050       0.4000            0.1625
## 6   741.6        0.1791         0.5249       0.5355            0.1741
##   sd symmetry sd fractal dimension
## 1      0.4601              0.11890
## 2      0.2750              0.08902
## 3      0.3613              0.08758
## 4      0.6638              0.17300
## 5      0.2364              0.07678
## 6      0.3985              0.12440
\end{verbatim}

\hfill\break
(2) Chuẩn hóa dữ liệu để đảm bảo rằng tất cả các biến có cùng quy mô.

\begin{Shaded}
\begin{Highlighting}[]
\NormalTok{df\_scaled }\OtherTok{\textless{}{-}} \FunctionTok{scale}\NormalTok{(df)}

\FunctionTok{apply}\NormalTok{(df\_scaled, }\DecValTok{2}\NormalTok{, mean)}
\end{Highlighting}
\end{Shaded}

\begin{verbatim}
##               diagnosis             mean radius            mean texture 
##            7.061834e-17           -1.379776e-16            6.157678e-17 
##          mean perimeter               mean area         mean smoothness 
##           -1.190559e-16            1.219971e-16            1.618800e-16 
##        mean compactness          mean concavity     mean concave points 
##           -7.605173e-17            3.929080e-17           -5.337457e-17 
##           mean symmetry  mean fractal dimension            worst radius 
##            1.679646e-16            4.810260e-16            3.721164e-17 
##           worst texture         worst perimeter              worst area 
##           -1.065855e-16            1.074950e-16            2.873226e-17 
##        worst smoothness       worst compactness         worst concavity 
##            1.290824e-16           -1.610869e-17           -5.978051e-17 
##    worst concave points          worst symmetry worst fractal dimension 
##            7.922107e-18            9.126277e-17            1.564424e-18 
##               sd radius              sd texture            sd perimeter 
##           -8.483720e-17            9.038321e-18            5.661517e-17 
##                 sd area           sd smoothness          sd compactness 
##            6.734248e-18           -2.269302e-16            1.261676e-17 
##            sd concavity       sd concave points             sd symmetry 
##            9.174370e-17            1.166365e-17            2.760417e-16 
##    sd fractal dimension 
##            2.002241e-16
\end{verbatim}

\begin{Shaded}
\begin{Highlighting}[]
\FunctionTok{apply}\NormalTok{(df\_scaled, }\DecValTok{2}\NormalTok{, sd)}
\end{Highlighting}
\end{Shaded}

\begin{verbatim}
##               diagnosis             mean radius            mean texture 
##                       1                       1                       1 
##          mean perimeter               mean area         mean smoothness 
##                       1                       1                       1 
##        mean compactness          mean concavity     mean concave points 
##                       1                       1                       1 
##           mean symmetry  mean fractal dimension            worst radius 
##                       1                       1                       1 
##           worst texture         worst perimeter              worst area 
##                       1                       1                       1 
##        worst smoothness       worst compactness         worst concavity 
##                       1                       1                       1 
##    worst concave points          worst symmetry worst fractal dimension 
##                       1                       1                       1 
##               sd radius              sd texture            sd perimeter 
##                       1                       1                       1 
##                 sd area           sd smoothness          sd compactness 
##                       1                       1                       1 
##            sd concavity       sd concave points             sd symmetry 
##                       1                       1                       1 
##    sd fractal dimension 
##                       1
\end{verbatim}

\hfill\break
(3) Tính toán ma trận hiệp phương sai và tìm giá trị riêng, vector
riêng.

\begin{Shaded}
\begin{Highlighting}[]
\NormalTok{cov\_matrix }\OtherTok{\textless{}{-}} \FunctionTok{cov}\NormalTok{(df\_scaled)}

\NormalTok{eigen\_result }\OtherTok{\textless{}{-}} \FunctionTok{eigen}\NormalTok{(cov\_matrix)}

\NormalTok{eigen\_values }\OtherTok{\textless{}{-}}\NormalTok{ eigen\_result}\SpecialCharTok{$}\NormalTok{values}
\NormalTok{eigen\_vectors }\OtherTok{\textless{}{-}}\NormalTok{ eigen\_result}\SpecialCharTok{$}\NormalTok{vectors}

\NormalTok{eigen\_values}
\end{Highlighting}
\end{Shaded}

\begin{verbatim}
##  [1] 1.391777e+01 5.726352e+00 2.846849e+00 1.998363e+00 1.659079e+00
##  [6] 1.207508e+00 6.847192e-01 4.848554e-01 4.168948e-01 3.508936e-01
## [11] 3.048554e-01 2.909858e-01 2.610104e-01 2.122756e-01 1.485762e-01
## [16] 8.816256e-02 7.986017e-02 5.903533e-02 5.146191e-02 4.821595e-02
## [21] 3.070628e-02 2.952665e-02 2.734467e-02 2.428359e-02 1.755221e-02
## [26] 1.547334e-02 8.056144e-03 6.869732e-03 1.588467e-03 7.416788e-04
## [31] 1.326347e-04
\end{verbatim}

\hfill\break
(4) Thực hiện PCA và phân tích kết quả. Sau đó, vẽ các biểu đồ liên quan
đến PCA, bao gồm biplot, biểu đồ tỷ lệ phương sai giải thích và biểu đồ
phân tán các thành phần chính.

\begin{Shaded}
\begin{Highlighting}[]
\NormalTok{pca }\OtherTok{\textless{}{-}} \FunctionTok{princomp}\NormalTok{(df\_scaled)}

\FunctionTok{summary}\NormalTok{(pca)}
\end{Highlighting}
\end{Shaded}

\begin{verbatim}
## Importance of components:
##                           Comp.1    Comp.2     Comp.3     Comp.4     Comp.5
## Standard deviation     3.7273732 2.3908760 1.68577763 1.41239197 1.28691989
## Proportion of Variance 0.4489604 0.1847210 0.09183385 0.06446333 0.05351866
## Cumulative Proportion  0.4489604 0.6336814 0.72551525 0.78997857 0.84349724
##                            Comp.6     Comp.7    Comp.8     Comp.9    Comp.10
## Standard deviation     1.09790061 0.82675013 0.6957035 0.64510631 0.59184200
## Proportion of Variance 0.03895187 0.02208771 0.0156405 0.01344822 0.01131915
## Cumulative Proportion  0.88244910 0.90453682 0.9201773 0.93362554 0.94494469
##                            Comp.11    Comp.12     Comp.13     Comp.14
## Standard deviation     0.551651721 0.53895681 0.510442654 0.460328668
## Proportion of Variance 0.009834045 0.00938664 0.008419691 0.006847598
## Cumulative Proportion  0.954778732 0.96416537 0.972585063 0.979432661
##                           Comp.15     Comp.16     Comp.17     Comp.18
## Standard deviation     0.38511695 0.296660781 0.282346987 0.242758270
## Proportion of Variance 0.00479278 0.002843954 0.002576135 0.001904366
## Cumulative Proportion  0.98422544 0.987069395 0.989645529 0.991549895
##                            Comp.19     Comp.20      Comp.21      Comp.22
## Standard deviation     0.226652742 0.219388265 0.1750780211 0.1716821445
## Proportion of Variance 0.001660062 0.001555353 0.0009905251 0.0009524726
## Cumulative Proportion  0.993209957 0.994765310 0.9957558349 0.9967083075
##                            Comp.23      Comp.24      Comp.25      Comp.26
## Standard deviation     0.165216858 0.1556949415 0.1323683001 0.1242825166
## Proportion of Variance 0.000882086 0.0007833417 0.0005662005 0.0004991399
## Cumulative Proportion  0.997590394 0.9983737352 0.9989399357 0.9994390756
##                             Comp.27      Comp.28      Comp.29      Comp.30
## Standard deviation     0.0896771168 0.0828109811 3.982054e-02 2.720984e-02
## Proportion of Variance 0.0002598756 0.0002216043 5.124088e-05 2.392512e-05
## Cumulative Proportion  0.9996989512 0.9999205555 9.999718e-01 9.999957e-01
##                             Comp.31
## Standard deviation     1.150659e-02
## Proportion of Variance 4.278539e-06
## Cumulative Proportion  1.000000e+00
\end{verbatim}

\begin{Shaded}
\begin{Highlighting}[]
\NormalTok{var }\OtherTok{\textless{}{-}}\NormalTok{ pca}\SpecialCharTok{$}\NormalTok{sdev}\SpecialCharTok{\^{}}\DecValTok{2} \SpecialCharTok{/} \FunctionTok{sum}\NormalTok{(pca}\SpecialCharTok{$}\NormalTok{sdev}\SpecialCharTok{\^{}}\DecValTok{2}\NormalTok{)}

\FunctionTok{plot}\NormalTok{(var,}
  \AttributeTok{col =} \StringTok{"red"}\NormalTok{,}
  \AttributeTok{cex =} \DecValTok{2}\NormalTok{,}
  \AttributeTok{lwd =} \DecValTok{2}\NormalTok{,}
  \AttributeTok{type =} \StringTok{"b"}\NormalTok{,}
  \AttributeTok{xlab =} \StringTok{"Thành phần chính"}\NormalTok{,}
  \AttributeTok{ylab =} \StringTok{"Tỉ lệ"}\NormalTok{,}
  \AttributeTok{main =} \StringTok{"Tỉ lệ các thành phần (phân tích trên bộ dữ liệu)"}
\NormalTok{)}
\end{Highlighting}
\end{Shaded}

\pandocbounded{\includegraphics[keepaspectratio]{W6_files/figure-latex/p2.4-1.pdf}}

\begin{Shaded}
\begin{Highlighting}[]
\FunctionTok{biplot}\NormalTok{(pca,}
  \AttributeTok{col =} \FunctionTok{c}\NormalTok{(}\DecValTok{2}\NormalTok{, }\DecValTok{4}\NormalTok{),}
  \AttributeTok{cex =} \FunctionTok{c}\NormalTok{(.}\DecValTok{5}\NormalTok{, }\FloatTok{1.25}\NormalTok{),}
  \AttributeTok{xlab =} \StringTok{"Thành phần chính thứ nhất"}\NormalTok{,}
  \AttributeTok{ylab =} \StringTok{"Thành phần chính thứ hai"}\NormalTok{,}
  \AttributeTok{main =} \StringTok{"Biểu đồ hai thành phần (sử dụng ma trận tương quan)"}
\NormalTok{)}
\end{Highlighting}
\end{Shaded}

\pandocbounded{\includegraphics[keepaspectratio]{W6_files/figure-latex/p2.4-2.pdf}}

\hfill\break
(5) Viết nhận xét về kết quả phân tích và ý nghĩa của các thành phần
chính.

Thành phần chính thứ nhất giải thích khoảng 45\% dữ liệu.

Thành phần chính thứ hai giải thích khoảng 18\% dữ liệu.

Thành phần chính thứ ba giải thích khoảng 9\% dữ liệu.

\ldots{}

7 thành phần chính đầu tiên giải thích khoảng 90\% dữ liệu.

11 thành phần chính đầu tiên giải thích khoảng 95\% dữ liệu.

17 thành phần chính đầu tiên giải thích khoảng 99\% dữ liệu.

\end{document}
